\section{Introduction}
\label{sec:intro}

Vietnam has one of the fastest-growing social-media user bases in Southeast Asia,
and the volume of Vietnamese text platforms must moderate keeps rising.  Flagging
hateful and offensive comments automatically is therefore a practical prerequisite
for healthier communities, yet Vietnamese remains low-resource and its social-media
style is hard (code-switching, ``teencode,'' informal spellings, evasive
substitutions, dropped diacritics).

The release of \emph{UIT-ViHSD} \cite{luu2021vihsd}, a human-annotated corpus of
roughly 33k Vietnamese comments labelled \textsc{clean}, \textsc{offensive} or
\textsc{hate}, made supervised transformer fine-tuning a realistic default and a
standard benchmark for monolingual encoders such as PhoBERT \cite{nguyen2020phobert}
and multilingual ones such as mBERT and XLM-R
\cite{conneau2020xlmr,nguyen2022smtce,hoang2023vihos}.  Multilingual DeBERTa-v3
variants have become popular because they combine a large multilingual pre-training
corpus with disentangled attention \cite{he2021deberta,he2023debertav3}.

In parallel, many \emph{optimisation} tricks are routinely applied to imbalanced
text classification: a richer classification head, discriminative (layer-wise)
learning rates, multi-stage fine-tuning, class-weighted or focal losses
\cite{lin2017focal}, curriculum learning \cite{bengio2009curriculum}, and per-class
decision thresholds \cite{saito2015precision}.  In the literature these are usually
bundled together and reported as a single end-to-end improvement, which leaves a
practical question unanswered: \emph{which of these decisions actually moves the
number on Vietnamese hate speech detection, and by how much?}

This paper answers that with a \emph{controlled ablation}: six systems sharing one
backbone (\texttt{microsoft/mdeberta-v3-base}), one fingerprinted split, one frozen
label mapping, one tokenizer configuration and one seed, differing \emph{only} in
the optimisation recipe and each scored once on the held-out test set.  The chain
(Sect.~\ref{sec:method}) runs from a plain baseline, to a multi-view pooling head,
to a variant adding discriminative rates and class-weighted two-stage training, to a
three-stage curriculum.

Our contributions are:

\begin{itemize}
  \item A comparison of six optimisation recipes for three-class Vietnamese hate
        speech detection under one fixed split and seed, reported with split
        fingerprints, per-class metrics, bootstrap intervals and paired significance
        tests so it can be reproduced rather than re-argued.
  \item A multi-view classification head (first-token, masked-mean and masked-max
        pooling) attaining the best macro-F1 at essentially the same parameter
        count as the standard head.
  \item An empirical demonstration that the recipes are not uniformly better: no
        pairwise macro-F1 difference is significant (a failure to detect, not
        proof of equivalence), the most accurate system has one of the lowest
        macro-F1 scores, a three-stage curriculum never wins validation
        selection, and a validation-searched threshold collapses to the trivial
        $0.5$---negative results we discuss for applied moderation.
\end{itemize}
